
\documentclass[12pt]{article}
\usepackage{amsmath, amssymb, amsfonts}
\usepackage{geometry}
\usepackage{hyperref}
\geometry{margin=1in}

\title{Unified Bead-Lattice Time-Shell Quantum Gravity:\\
A Discrete Framework for Emergent Spacetime, Gravitons, and Cosmology}
\author{Anonymous Submission}
\date{}

\begin{document}
\maketitle

\begin{abstract}
We present a unified theoretical framework connecting Planck-scale discreteness,
quantum gravity, classical general relativity, entanglement-induced topology,
and cosmology. The model draws on and extends ideas from discrete quantum 
gravity~\cite{Regge1961}, holography~\cite{Susskind1995}, ER=EPR 
entanglement geometry~\cite{Maldacena2013}, quantum foam~\cite{Wheeler1957}, 
and induced gravity mechanisms~\cite{Sakharov1967}. A detailed set of 
appendices provides technical derivations.
\end{abstract}

\section{Introduction}
Our model builds upon the view that spacetime geometry emerges from 
microscopic degrees of freedom~\cite{Jacobson1995, Verlinde2011}. 
We introduce a bead-lattice structure on a dynamic time shell, from which 
gravitons, curvature, wormholes, and cosmological evolution emerge.

\section{Planck-Scale Bead Structure}
Discrete spacetime approaches such as Regge calculus~\cite{Regge1961} 
and loop quantum gravity~\cite{Rovelli1998} motivate the bead model. 
Each bead carries geometric and energetic variables whose collective dynamics
define curvature.

\section{Time-Shell Geometry}
The time-shell structure resonates with holographic principles where 
2D surfaces encode 3D bulk information~\cite{Susskind1995, Bousso2002}.

\section{Unified Lagrangian}
The action's bulk components reflect standard quantum field 
theory in curved space~\cite{BirrellDavies}, while the shell structure 
extends induced gravity ideas~\cite{Sakharov1967}.

\section{Linearized Gravity and TT Modes}
The emergence of massless gravitons follows 
Fierz–Pauli theory~\cite{FierzPauli1939} and modern effective field theory 
approaches~\cite{Donoghue1994}.

\section{ER Bridges and Entanglement Geometry}
ER=EPR~\cite{Maldacena2013} indicates a deep relationship between 
entanglement and wormhole geometry. Our bead adjacency model extends this to 
Planck-scale discreteness.

\section{Continuum Limit and GR}
Classical Einstein gravity emerges following arguments similar to 
Sakharov~\cite{Sakharov1967} and entropic gravity formulations 
like~\cite{Jacobson1995, Verlinde2011}.

\section{Cosmology}
Our time-shell cosmology aligns conceptually with holographic cosmic 
models~\cite{Padmanabhan2010}.

\appendix
\section{Appendix A: TT Derivations}
See~\cite{FierzPauli1939, Donoghue1994} for classical references on 
linearized gravity.

\section{Appendix B: Discrete Hamiltonian}
Related discrete formulations appear in~\cite{Regge1961}.

\section{Appendix C: Continuum Limit}
Emergence of GR via coarse-graining has connections to 
entanglement-gravity arguments~\cite{Jacobson1995}.

\section{Appendix D: ER-Bridge Geometry}
The relation between entanglement entropy and spacetime geometry is 
summarized in~\cite{Maldacena2013}.

\section{Appendix E: Cosmological Cycles}
Connections to holographic cosmology appear in~\cite{Padmanabhan2010}.

\bibliographystyle{unsrt}
\bibliography{peer_review_refs}

\end{document}
